We will begin our environment setup by building Clangd. Then, we will install VS Code, set up the Clangd extension, and configure Clangd within it.

\mySubsubsection{8.3.1.}{Clangd build}

It’s worth building Clangd in release mode, as we did for LLDB in Section 1.3.3, The LLVM debugger, its build, and usage. This is because performance is crucial in IDEs. For instance, Clangd requires building an AST to provide code navigation functionality. If a user modifies a document, the document should be rebuilt, and the navigation functionality will not be available until the rebuild process is completed. This can result in delays in IDE responses. To prevent IDE slowness in responses, we should ensure that Clangd is built with all the required optimizations. You can use the following project configuration command:

\begin{shell}
cmake -G Ninja -DCMAKE_BUILD_TYPE=Release -DCMAKE_INSTALL_PREFIX=../install -DLLVM_TARGETS_TO_BUILD="X86" -DLLVM_ENABLE_PROJECTS="clang;clang-tools-extra" ../llvm
\end{shell}

\begin{center}
Figure 8.3: Release configuration for Clangd build
\end{center}

The command has to be run from the release folder that we created in Section 1.3.3, The LLVM debugger, its build, and usage, as shown in Figure 1.8. As you can see, we have enabled two projects in Figure 8.3: clang and clang-tools-extra .

You can use the following command to build and install Clangd:

\begin{shell}
$ ninja install-clangd -j $(nproc)
\end{shell}

This command will utilize the maximum available threads on the system and install the binary into the folder specified in our CMake command in Figure 8.3, which is the install folder under the LLVM source tree.

After building the Clangd binary, our next step will include installing VS Code and configuring it to work with Clangd.

\mySubsubsection{8.3.2.}{VS Code installation and setup}

You can download and install VS Code from the VS Code website: \url{https://code.visualstudio.com/download}.

The first step after running VS Code is to install the Clangd extension. An open source extension is available to work with Clangd via LSP. The extension’s source code can be found on GitHub: https://github.com/clangd/vscode-clangd. However, we can easily install the latest version of the extension directly from within VS Code.

To do this, press Ctrl+Shift+X (or \includegraphics[width=0.02\textwidth]{content/part2/chapter8/images/3.png}+Shift+X for macOS) to open the extensions panel. Search for Clangd and click the Install button.

\myGraphic{0.5}{content/part2/chapter8/images/4.png}{Figure 8.4: Installing the Clangd extension}

After installing the extension, we need to set it up. The main step is to specify the path to the Clangd executable.

You can access this setting via the File — Preferences — Settings menu or by pressing Ctrl + , (or\includegraphics[width=0.02\textwidth]{content/part2/chapter8/images/3.png} +, for macOS), as shown in the following screenshot:

\myGraphic{0.5}{content/part2/chapter8/images/5.png}{Figure 8.5: Setting up the Clangd extension}

As shown in Figure 8.5, we have configured the Clangd path to be /home/ivanmurashko /clangbook/llvm-project/install/bin/clangd . This path was used during the installation of the Clangd binary in Section 8.3.1, Clangd build.

You can open your favorite C/C++ source file and try to navigate through it. For instance, you can search for a definition for a token, switch between a source and a header, and so on. In our next example, we will investigate how navigation, and especially go-to definition, works through LSP.

\begin{myNotic}{Important note}
Our setup works only for simple projects that do not require special compilation flags. If your project requires special configuration to build, then you have to use a generated compile\_commands.json file that should be placed at the root of your project. This file should contain a Compilation Database (CDB) in JSON format, specifying compilation flags for each file in your project. For more information about the setup, please refer to Figure 9.5, Clangd Setup for Large Projects.
\end{myNotic}

With the required components installed, we are now ready for an LSP demo where we will emulate typical development activities in an IDE (open and modify a document, jump to a token definition, etc.) and explore how it’s represented via LSP.










































































